% Week 6 section of the synthesis exercise sheet, included by Exercises.tex.
\section{Week 6 --- Densities, Girsanov, and Measure Changes}

Throughout, all density processes are assumed adapted to the stated
filtration. Whenever a stochastic exponential is used, distinguish the
formal drift calculation from the condition that makes it a true martingale.

\subsection*{Part I --- What a Density Really Encodes}

\begin{exercise}[Easy to medium --- Reversing a change of measure]
Let \(Z>0\) satisfy \(\E^{\Pb}[Z]=1\), and define
\(\dd\Qbb=Z\,\dd\Pb\).
\begin{exparts}
  \item Prove that
  \[
  \frac{\dd\Pb}{\dd\Qbb}=\frac1Z
  \qquad\text{and}\qquad
  \E^{\Qbb}\!\left[\frac1Z\right]=1.
  \]
  \item Show that for every integrable \(Y\),
  \[
  \E^{\Pb}[Y]=\E^{\Qbb}\!\left[\frac{Y}{Z}\right].
  \]
  Explain why the random variable itself has not changed.
  \item Use Jensen's inequality to prove
  \[
  \E^{\Pb}[\log Z]\leq0.
  \]
  Assuming the relevant integrability, prove also
  \(\E^{\Qbb}[\log Z]\geq0\).
  \item Characterize the equality case in both inequalities. What does it
  say about the relation between \(\Pb\) and \(\Qbb\)?
\end{exparts}
\end{exercise}

\begin{exercise}[Medium --- A density process in a two-period tree]
Let
\[
\Omega=\{UU,UD,DU,DD\},
\]
where the first letter is observed at time \(1\) and the second at time \(2\).
Under \(\Pb\), all four paths have probability \(1/4\). Under \(\Qbb\), their
probabilities in the displayed order are
\[
\frac38,\qquad\frac18,\qquad\frac18,\qquad\frac38.
\]
Let \(\F_1\) reveal only the first letter, and let
\(Z_2=\dd\Qbb/\dd\Pb\).
\begin{exparts}
  \item Compute \(Z_2\) on every terminal path and verify
  \(\E^{\Pb}[Z_2]=1\).
  \item Compute \(Z_1=\E^{\Pb}[Z_2\mid\F_1]\). Explain why \(Z_1\) is the
  density of the time-\(1\) information under \(\Qbb\) relative to \(\Pb\).
  \item Compute under \(\Qbb\) the conditional probability that the second
  letter is \(U\), first given that the first letter is \(U\), then given
  that it is \(D\).
  \item Recover both answers from Bayes' formula
  \[
  \Qbb(A\mid\F_1)
  =\frac{\E^{\Pb}[Z_2\1_A\mid\F_1]}{Z_1}.
  \]
  Does the second step remain independent of the first under \(\Qbb\)?
\end{exparts}
\end{exercise}

\subsection*{Part II --- Tilting Beyond the Classroom Examples}

\begin{exercise}[Medium --- Tilting a correlated Gaussian vector]
Under \(\Pb\), let \((X,Y)\) be centered jointly Gaussian with
\[
\operatorname{Var}(X)=\sigma_X^2,\qquad
\operatorname{Var}(Y)=\sigma_Y^2,\qquad
\operatorname{Cov}(X,Y)=\rho\sigma_X\sigma_Y.
\]
For \(\theta\in\R\), define
\[
Z=\exp\!\left(\theta X-\frac12\theta^2\sigma_X^2\right),
\qquad \dd\Qbb=Z\,\dd\Pb.
\]
\begin{exparts}
  \item Verify that \(\E^{\Pb}[Z]=1\).
  \item Use the joint moment-generating function to compute
  \(\E^{\Qbb}[e^{uX+vY}]\).
  \item Identify the mean vector and covariance matrix of \((X,Y)\) under
  \(\Qbb\). Explain why tilting with \(X\) also moves the mean of \(Y\)
  when \(\rho\neq0\).
  \item Determine for which \(\rho\) the law of \(Y\) is unchanged by this
  measure change.
\end{exparts}
\end{exercise}

\begin{exercise}[Medium --- An adapted tilt made safe by stopping]
Let
\[
\tau_a=\inf\{t\geq0:|W_t|\geq a\}\wedge T,
\qquad
\phi_t=\theta W_t\1_{\{t\leq\tau_a\}},
\]
where \(a,T>0\) and \(\theta\in\R\). Define
\[
Z_t=\exp\!\left(
\int_0^t\phi_s\dd W_s-\frac12\int_0^t\phi_s^2\dd s
\right).
\]
\begin{exparts}
  \item Show that \(|\phi_t|\leq|\theta|a\), and verify Novikov's condition.
  Conclude that \(Z\) is a true martingale.
  \item Define \(\Qbb\) by \(Z_T\). Identify the \(\Qbb\)-Brownian motion
  \(B\) and write the \(\Qbb\)-dynamics of \(W\).
  \item Describe the drift of \(W\) before and after \(\tau_a\). For
  \(\theta>0\), does the pre-stopping drift push the process toward or away
  from zero?
  \item Explain why stopping is doing genuine mathematical work here rather
  than merely shortening the notation.
\end{exparts}
\end{exercise}

\subsection*{Part III --- Designing and Testing a Measure Change}

\begin{exercise}[Medium --- Changing one mean-reversion model into another]
Under \(\Pb\), suppose
\[
\dd X_t=\kappa(m-X_t)\dd t+\sigma\dd W_t^{\Pb},
\qquad \kappa,\sigma>0.
\]
We would like \(X\) to satisfy under \(\Qbb\)
\[
\dd X_t=\widetilde\kappa(\widetilde m-X_t)\dd t
+\sigma\dd W_t^{\Qbb},
\qquad \widetilde\kappa>0.
\]
\begin{exparts}
  \item Find the candidate process \(\lambda_t\) in
  \(W_t^{\Qbb}=W_t^{\Pb}+\int_0^t\lambda_s\dd s\).
  \item Write the corresponding candidate density process, with the correct
  sign.
  \item Explain why the volatility must remain \(\sigma\). Could an
  equivalent Girsanov change turn it into a different constant
  \(\widetilde\sigma\neq\sigma\)?
  \item The candidate \(\lambda_t\) is generally unbounded. Explain why the
  algebraic drift match does not by itself prove that \(\Qbb\) exists, and
  describe how localization by
  \(\tau_n=\inf\{t:|X_t|\geq n\}\wedge T\) helps.
\end{exparts}
\end{exercise}

\begin{exercise}[Medium --- Can one Brownian shift price two stocks?]
Two traded stocks are driven by the same Brownian motion:
\[
\frac{\dd S_t^1}{S_t^1}=\mu_1\dd t+\sigma_1\dd W_t,
\qquad
\frac{\dd S_t^2}{S_t^2}=\mu_2\dd t+\sigma_2\dd W_t,
\]
where \(\sigma_1,\sigma_2\neq0\), and the short rate is \(r\).
\begin{exparts}
  \item If \(W_t^{\Qbb}=W_t+\lambda t\), derive the condition on
  \(\lambda\) that makes each discounted stock driftless.
  \item Show that a common equivalent martingale measure of this form can
  exist only if
  \[
  \frac{\mu_1-r}{\sigma_1}
  =
  \frac{\mu_2-r}{\sigma_2}.
  \]
  Interpret this equality economically.
  \item When the equality holds, identify the density process and explain
  why the Brownian risk premium is unique.
  \item Suppose the equality fails. Find dollar exposures $x_1,x_2$ in the
  two stocks such that $x_1\sigma_1+x_2\sigma_2=0$, while
  $x_1(\mu_1-r)+x_2(\mu_2-r)\neq0$. Explain why this is an instantaneous
  arbitrage signal and how the sign of the position can be chosen.
\end{exparts}
\end{exercise}
